\section{从继统之争到富国强兵}

1063年仁宗去世，英宗继位。仁宗无存活皇子的继承安排，需要宗室、后族、宰执和礼官共同承认；《宋史》两朝本纪和《长编》嘉祐八年条可确认这一转换。1065年的濮议随即把英宗生父濮王应获何种追尊推到朝廷中心。它不是可以被还原成纯粹礼学或纯粹党争的一次争吵：宗法名分怎样写入国家礼制，会决定谁有资格解释皇统，也会重新结连台谏和宰执。

1067年神宗即位，财政、兵役和边防的未决问题进入新阶段。次年王安石入朝，神宗与他讨论变法方向；《宋史》神宗本纪、王安石传和《长编》提供了人事与论议的骨架。把其后的全部政策都写成一人设计、全国照办的蓝图，会遮蔽皇帝、制置机构、地方官和胥吏各自的作用。新法首先是一套试图把财政来源、基层劳役与军事压力重新接合的制度实验。

1069年青苗法颁行，官府在青黄不接时以贷款介入借贷；1070年前后免役法在部分地区推进，以雇募和免役钱改造差役；1071年保甲法开始以户籍编组治安、训练与互保。《宋史》食货志、兵志、《宋会要辑稿》与《长编》能分别确认这些制度文本和推进节点。它们不能证明同一利率、同一征收方法或同一训练密度覆盖全国。对小民而言，政策的实际面貌取决于户等认定、地方官考成、书手催科和市场中的货币；正是在这些环节，救济、加派、逃避与争讼可能并存。

\section{新法走向边境，也走进文字}

1072年市易法设市易务，以官本介入城市商货和信贷。它说明改革并未只瞄准乡村：开封等城市已是大宗货物、借贷与官府需求交会之处。会要和食货志所记的机构、资金并不能单独证明物价稳定或商人盈亏；还须由区域性市场材料检验。制度争论的核心也不只是“官”与“商”的抽象对立，而在于谁掌握货物信息、何时催收本息、谁承担价格波动。

1074年旱灾、政策批评和人事冲突交织，王安石首次罢相。灾异在政治语言中常被用作责难政府的证据，不能据此断言旱灾直接造成政策逆转。此后改革并未终止，而是更依赖神宗与不同官僚网络的支持。1079年苏轼因诗文被控讥讽朝政，乌台诗案显示台谏、法司和文人如何以诗句、奏议和“结党”语言划定忠诚边界；《宋史》苏轼传与《长编》可证案件与处分，却不容把每句诗的政治含义写成唯一解释。

\subsection{交读窗口：黄州的水面与乌台案的文书}

元丰五年（一〇八二）秋，仍在黄州的苏轼写《前赤壁赋》。这是以舟游、问答和主客转换铺写的散赋，不是申诉状或判牍；“寄蜉蝣于天地，渺沧海之一粟”把遭贬者置于江水与时间的尺度中。它之所以能照亮乌台诗案之后的政治文化，不在于替案件供给一条新事实，而在于让人看见一位被处分的官员如何把失意、友游与经典修辞转换为可公开流传的文字。

案件的日期、牵连的奏议与处分，仍应以《续资治通鉴长编》元丰二年条、《宋史》〈苏轼传〉及可追查的《宋会要辑稿·刑法》条目为脊梁；赋中“客”与“苏子”的声音也不能逐一还原成法司、台谏或黄州居民的意见。钱穆所关心的制度与士人之间如何彼此制约，在此可转化为一个具体问题：案件后的言说空间由谁界定？吕思勉式的社会史追问则要求把诗文名士与催科、保甲、城寨中的人分开观察。地方墓志、题名和文集的保存尤偏向识字精英，故《前赤壁赋》可以说明一段文学姿态，不能量度新法执行，亦不能代表所有被政令卷入者的感受。

边疆的账目同样没有消失。1075年李朝军队攻入邕、钦、廉州，1076--1077年郭逵率军南征至富良江一线后撤。宋、越材料都承认战争，却在战果、疫病与损失上保留各自叙事位置；热带道路和补给的限制，比单一的勇怯判断更能解释未能长期占领。1080年永乐城失守，1081年五路伐夏未获预期战果，城寨的选址、供水、运输和多路协同再度暴露为军事问题。1082年元丰改制重整官制名目与机构关系，意在处理权责交错；复古的官名不等于实际职能立刻稳定。新法的代价与边防的代价，都仍由中枢的文书、财政和地方执行来承担。
